\section{The Pearl Tiara}

A pearl tiara has a design and a structure. The pearls are arranged in a certain way, usually by size. It is built around precious metals and may include other gems. Although the quality of each individual pearl is important, the piece of jewelry achieves its beauty in its totality, not from any pearl in particular. This is opposed to a bowl of pearls, in which each pearl has no relation to its neighbor, that is, it is a heap. In a world that considers a bowl of pearls to be a work of art, there are obviously many who cannot discern the difference.

Blogs, too, can be read in two different ways. Careless readers will regard the blog posts like pearls in a bowl, with nothing but casual relationships among the various posts. They cannot see that Gornahoor is architectonic, posts are related to each other, themes appear and reappear with variations, and that nothing is posted without a reason nor without a place in the totality. Although this design is not visible to the senses as a pearl tiara is, it should be visible to the intellect.

Besides the text itself, there is an art. The coat of arms, the motto, the colours, etc., are all planned for a certain effect. The vocabulary is deliberately chosen so that only those of a certain intelligence and education can understand it, yet any difficulty is not from intentional obfuscation.

\paragraph{Tone Deaf to Life}


For example, one of our first posts, Tone Deaf to Life\footnote{\url{https://gornahoor.net/?p=12}}, published in 2007, illustrates this. On the one hand, it starts with the notion of the consciousness of the True Self (or Real I). It ends with the challenge to see one's life in its totality, i.e., like a tiara, rather than a collection of random events like a bowl.

Seven years later, we come back to this, but a little deeper. Instead of isolated quotes, we provide entire, previously untranslated, texts from Julius Evola on the same them of the consciousness of the true self. Furthermore, beyond each individual post, there is an implicit, if not explicit, relationship among the posts. In other words, Gornahoor as a blog tries to be, in microcosm, an exemplar of the macrocosm it describes.

And not just rhetorically, since that is the recent theme. We are now also providing the opportunity and the venue\footnote{\url{https://gornahoor.net/?p=7205}} for some among you to actually develop the consciousness of the true self as a reality, as being, and not just as text, rhetoric, or thought.

\paragraph{The Attraction of Rhetoric}


So when I read of a suggestion to do a comparative study of Evola, Michelstaedter, and Heidegger, I should feel frustrated, except that I understand the human, all too human, condition too well. It suggests rhetoric in a post that is anti-rhetoric. Such a study is suitable for a professor who needs to publish in between his various conferences and seductions of his graduate assistants.

To look at things in their totality: it starts with Evola's intention to discuss five earlier thinkers. In no way, did he write a comparative study of those five men. Rather, he had a different purpose. He saw his philosophy as the fulfillment of what those men may have started, but never completed. Specifically, he wrote:

\begin{quotex} 
Magical idealism carries forward by way of logical continuity and integrates the most advanced positions that modern Western speculation has achieved
\end{quotex}


How could anyone have missed that? It is the opposite of rhetoric, Heidegger's \emph{gerede} or idle chitchat. Instead, in an act of possession, Evola incorporates the best of those writers into his own scheme, as he claims. He brings out the logical consequences of the earlier writers, consequences that they themselves may have overlooked. He integrates their positions into a larger whole.

Is that vanity, or self-aggrandizement, or confidence? Or perhaps the words of one who knows?

Can anyone take the next logical step? Why would we would we bother to translate and publish such an obscure text? Evola's text is the object and Gornahoor is the frame. Next is to see Gornahoor as the object and then determine the frame.

Perhaps you have guessed by now. Evola announced a project and a method. Since we take such writings seriously and not as mere rhetoric used to provide intellectual entertainment, we take possession of that project and method. Hence, Gornahoor is best understood as the historical fulfillment of the most advanced positions that Western tradition has achieved.

This has been done by bringing out the logical consequences of previous writers and by integrating their positions into a larger whole. The list of writers referenced is large, but not that large. Does this help anyone seek the tiara rather than the bowl?

So, let us go back to the original proposal. An intelligent man will write about Evola, Michelstaedter, and Heidegger, not to exalt them, compare them, or contrast them, but to overcome them, fulfill them, and integrate them into a larger whole.


\flrightit{Posted on 2014-06-06 by Cologero}

\begin{center}* * *\end{center}

\begin{footnotesize}
\begin{sffamily}
\texttt{Michael on 2023-06-06 at 09:48 said:}

It is important for reader's to understand the goal of this post. To truly understand and incorporate these thinkers is to sit with them and not try to resolve them with each other, but to resolve them within oneself. Producing something from that point brings the current to the surface.

\end{sffamily}
\end{footnotesize}

